%% =====================================================================
%%  T-08-01  Do Not Offend the Wrong Person
%%  -------------------------------------------------------------------
%%  One idea per file. Core is mandatory. Slots in canonical order.
%%  Max 700 words. No layout commands. No filler. New source material
%%  for this entry goes in sources/<BOOK>/T-08-01.tex -- never by
%%  rewriting this file (docs/RULES.md R0.1).
%% =====================================================================
%% @kind: topic
%% @id: T-08-01
%% @title: Do Not Offend the Wrong Person
%% @subtitle: The cost of engagement depends entirely on who you engage
%% @chapter: c08
%% @order: 10
%% @status: stable
%% @sources: B3
%% @updated: 2026-09-19
%% @allow-rewrite: slot migration to the seven-question format (docs/RULES.md section 2)

\topic{T-08-01}{Do Not Offend the Wrong Person}{The cost of engagement depends entirely on who you engage}


\begin{core}
Most losses in this domain come from choosing the counterpart badly rather than from executing badly. Some people respond to being outmanoeuvred by forgetting it; others spend the rest of their lives on it, and you cannot tell which by looking at how they behave toward you today.
\end{core}

\begin{mechanism}
The variance in response to a slight is enormous and is not visible from the front. Ordinary people absorb small defeats and move on; the dispositions described in T-03-01 do not, and their response is delayed, disproportionate and sustained. Because the response is delayed, the person who caused it usually has no idea it is coming and no idea what caused it --- the offence was minor, the interaction ended amiably, and the account was opened anyway. The practical consequence is that selection matters more than technique: an average move against a safe counterpart beats a brilliant move against a dangerous one.
\end{mechanism}

\begin{application}
  \item Assume variance. Do not infer response type from the current interaction.
  \item Gather history before engaging: how did this person handle the last person who crossed them?
  \item Keep the cost of any move below the threshold at which a delayed account is opened.
  \item Where you must oppose a dangerous counterpart, oppose the proposal rather than the person, and leave them a face-saving route.
  \item Where you have already offended one, settle it completely and early. Partial settlement is worse than none.
\end{application}

\begin{feedback}
  \item Disproportionate reaction to something minor, especially in public.
  \item Coldness at the time of an injury, followed by nothing --- and then, much later, by something.
  \item A long, detailed memory of how they were treated.
  \item A history of conflicts with former colleagues, partners or friends, all of which were the other party's fault.
  \item Warmth that arrives very quickly, before there is anything to be warm about.
\end{feedback}

\begin{failure}
It counsels caution to the point of paralysis if taken literally, since no counterpart's response type is fully observable in advance. It also has a moral cost: a policy of never offending the dangerous and freely offending the safe is a policy of predation on the people least able to defend themselves, and the source eventually prices that. The usable version is a screening rule --- gather history, keep costs low, settle early --- not a targeting rule.
\end{failure}

\begin{countermeasures}
  \item Ask who this person has fallen out with, and how it ended. The answer is available and is the single most predictive thing about them.
  \item Do not assume you are the exception. Special-exemption claims are covered in T-06-03.
  \item Keep your own exposure low with anyone whose response type you have not yet observed under conflict --- which is to say, with almost everyone, early on.
  \item If you are the person who has been offended, notice that delayed disproportionate response is expensive for you as well as for them, and settle or exit instead of accruing.
\end{countermeasures}

\sources{B3}
\seesources
\seealso{T-03-01, T-08-02, T-21-01}
